% --- Chapter 4 Lead-In ---
The unprotected limit study showed that instruction memory fails deterministically once faults persist in static storage. 
Resets and detection mechanisms provide no recovery because corrupted instructions are re-executed identically after every reboot. 
This behavior defines the lower bound of software reliability—where faults in code itself permanently prevent progress.  

To move beyond detection, the system must regenerate correct behavior from within the same faulty memory space. 
\textbf{Repair} achieves this by algorithmically reconstructing code or data from redundant or overlapping information, 
producing a more reliable version of the original content. 
Unlike reset-based recovery, repair operates in place, fusing imperfect data to restore an executable state.  

This chapter introduces repair as the first active recovery mechanism for static memory faults. 
It evaluates two representative models—\textbf{Majority-of-3} and \textbf{Reed–Solomon (255, 223)}—that combine information from multiple unreliable regions to reconstruct usable code. 
Both approaches operate entirely within memory, without reboot or external reprogramming, and aim to restore a coherent instruction stream by reinforcing correct content through redundancy.  

\section{Overview}
\label{sec:repair-overview}

The following sections develop a quantitative model of repair, derive analytical expressions for survivability, 
and validate these models through a limit study that measures how repair strength and redundancy influence system reliability. 
These results establish repair as a baseline mechanism for sustaining execution when persistent instruction corruption 
would otherwise render the system non-recoverable.